\documentclass{tufte-handout}

\usepackage{style}

\begin{document}

\tma{03} % TMA number

\begin{question}

\qpart

\qsubpart

\[ \left(\frac{-48}{229}\right) \]

\begin{align*}
\stext{\textup{Theorem 2.2(c), HB Chapter 10 p46}}
\left(\frac{-48}{229}\right) &= \left(\frac{-1}{229}\right) \cdot \left(\frac{16}{229}\right) \cdot \left(\frac{3}{229}\right)\\[8pt]
\stext{Calculating each term separately, using standard results}
\left(\frac{-1}{229}\right) &= 1\\[8pt]
\snote{\textup{Theorem 2.2(e), HB Chapter 10 p46}}\\
\left(\frac{16}{229}\right) &= 1\\[8pt]
\snote{\textup{Theorem 2.2(b), HB Chapter 10 p46}}\\
\left(\frac{3}{229}\right) &= 1\\[8pt]
\snote{\textup{Proposition 4.7, HB Chapter 10 p46}}\\
\stext{Combining these results}
\left(\frac{-48}{229}\right) &= 1 \cdot 1 \cdot 1 \\[8pt]
&= 1 \\[8pt]
\stext{So \( -48 \) is a quadratic residue modulo \( 229 \)}
\end{align*}

\vspace{2cm}

\qsubpart

\[ \left(\frac{151}{43}\right) \]

\begin{align*}
\left(\frac{151}{43}\right) &= \left(\frac{22}{43}\right) \\[8pt]
&= \left(\frac{2}{43}\right) \cdot \left(\frac{11}{43}\right) \\[8pt]
\snote{\textup{Theorem 2.2(c), HB Chapter 10 p46}}\\
\stext{Calculating each term separately}
\left(\frac{2}{43}\right) &= -1 \\[8pt]
\snote{\textup{Proposition 3.2, HB Chapter 10 p46}}\\
\left(\frac{11}{43}\right) &= -\left(\frac{43}{11}\right) \\[8pt]
\snote{\textup{Theorem 4.2, HB Chapter 10 p46}}\\
&= -\left(\frac{10}{11}\right) \\[8pt]
&= -\left(\frac{2}{11}\right) \cdot \left(\frac{5}{11}\right) \\[8pt]
\snote{\textup{Theorem 2.2(c), HB Chapter 10 p46}}\\
\stext{Calculate each term seperatly}
\left(\frac{2}{11}\right) &= -1 \\[8pt]
\snote{\textup{Proposition 3.2, HB Chapter 10 p46}}\\
\left(\frac{5}{11}\right) &= \left(\frac{11}{5}\right) \\[8pt]
\snote{\textup{Theorem 4.2, HB Chapter 10 p46}}\\
&= \left(\frac{1}{5}\right) \\[8pt]
\left(\frac{1}{5}\right) &= 1 \\[8pt]
\snote{\textup{Theorem 2.2(b), HB Chapter 10 p46}}\\
\stext{Combining these results}
\left(\frac{151}{43}\right) &= \left(\frac{22}{43}\right)=\left(\frac{2}{43}\right)\left(\frac{11}{43}\right)\\[8pt]
\left(\frac{2}{43}\right) &=  -1\\[8pt]
\stext{and}
\left(\frac{11}{43}\right) &= -\left(\frac{10}{11}\right) = -\left(\frac{2}{11}\right) \cdot \left(\frac{5}{11}\right)\\
\left(\frac{2}{11}\right) &= -1\\[8pt]
\left(\frac{5}{11}\right) &= 1\\[8pt]
\stext{So}
\left(\frac{11}{43}\right) &=  -\left(\frac{2}{11}\right) \cdot \left(\frac{5}{11}\right) = -((-1)(1)) = 1\\[8pt]
\stext{therefore}
\left(\frac{151}{43}\right) &= \left(\frac{2}{43}\right)\left(\frac{11}{43}\right) = (-1)(1) = -1\\[8pt]
\stext{So \( 151 \) is a quadratic non-residue modulo \( 43 \)}
\end{align*}

\end{question}

\end{question}

\end{document}